\chapter{Dzielniki elementarne i bloki pierwotne}
\label{ch:lin20-elementary}
\index{dzielnik elementarny@Dzielnik elementarny}
\index{blok pierwotny@Blok pierwotny}
\index{rozkład pierwotny@Rozkład pierwotny}
\index{klasyfikacja podobieństwa@Klasyfikacja podobieństwa}

\noindent\textbf{Cel rozdziału.}
Postać Frobeniusa z Rozdziału~\ref{ch:lin19-frobenius}
opisuje operator przez uporządkowany ciąg
czynników niezmienniczych. Każdy z nich
można rozłożyć nad ustalonym ciałem
na potęgi wielomianów nierozkładalnych.
Uzyskamy w ten sposób \emph{dzielniki elementarne}
i drugą postać blokową operatora.
Pokażemy, jak odzyskiwać je z wymiarów jąder
wielomianów operatora oraz jak łączą się
z blokami Jordana. Nie będziemy ponownie
dowodzić ogólnego rozkładu pierwotnego
z Rozdziału~\ref{ch:lin16-niezmiennicze}.

W całym rozdziale \(K\) oznacza dowolne ciało,
\(V\) skończenie wymiarową przestrzeń wektorową
nad \(K\), \(n=\dim_K V\), a
\(T\in\operatorname{End}_K(V)\).
Wielomiany nierozkładalne przyjmujemy moniczne.
Jeżeli \(n=0\), wszystkie listy bloków
i czynników są puste, a
\(\chi_T=m_T=1\).

\section{Wielomiany nierozkładalne i twierdzenie o resztach}
\label{sec:lin20-crt}

\begin{definicja}[Dzielniki elementarne]
\label{def:lin20-elementary}
Niech \(p\in K[t]\) będzie monicznym
wielomianem nierozkładalnym
o dodatnim stopniu, a \(e\ge1\)
liczbą całkowitą.
Wielomian \(p^e\) nazywamy
\emph{potęgą pierwotną}.
\emph{Dzielnikami elementarnymi}
operatora \(T\) nazywamy wszystkie
potęgi \(p^e\), z uwzględnieniem powtórzeń,
otrzymane z rozkładów na czynniki
nierozkładalne jego czynników niezmienniczych.
Odpowiadający wielomianowi \(p^e\)
blok \(C(p^e)\) nazywamy
\emph{blokiem pierwotnym}.
\end{definicja}

Każdy niezerowy wielomian moniczny
dodatniego stopnia jest iloczynem monicznych
wielomianów nierozkładalnych.
Istnienie takiego rozkładu uzasadnia
indukcja po stopniu: jeżeli wielomian
jest rozkładalny, oba jego nieodwracalne
czynniki mają mniejsze dodatnie stopnie.
Jednoznaczność z dokładnością do kolejności
wynika z lematu Euklidesa dla \(K[t]\):
jeżeli nierozkładalny \(p\) dzieli \(ab\),
to dzieli \(a\) lub \(b\).
Istotnie, gdy \(p\nmid a\),
algorytm Euklidesa daje
\(up+va=1\), a pomnożenie przez \(b\)
pokazuje \(p\mid b\).
Ponieważ czynniki normalizujemy do postaci
monicznej, nie pozostaje nieoznaczoność
związana z mnożeniem przez stałe.
Dla wielomianu \(1\) rozkład jest pusty.

\begin{twierdzenie}[Chińskie twierdzenie o resztach dla {\(K[t]\)}]
\label{thm:lin20-crt}
Niech \(f_1,\ldots,f_s\in K[t]\),
\(s\ge1\), będą monicznymi wielomianami
dodatniego stopnia, parami względnie pierwszymi.
Wówczas odwzorowanie
\[
 \Phi\colon K[t]/(f_1\cdots f_s)
 \longrightarrow\bigoplus_{i=1}^{s}K[t]/(f_i),
 \qquad
 [g]_{f_1\cdots f_s}\longmapsto
       ([g]_{f_1},\ldots,[g]_{f_s})
\]
jest izomorfizmem modułów nad \(K[t]\).
\end{twierdzenie}
\begin{proof}
Jeżeli dwa wielomiany różnią się
o wielokrotność \(f_1\cdots f_s\),
to różnią się także o wielokrotność
każdego \(f_i\). Dlatego \(\Phi\)
jest dobrze określone.
Wprost ze wzoru wynika jego liniowość
nad \(K[t]\).

Jeśli \(\Phi([g])=0\), to
\(f_i\mid g\) dla każdego \(i\).
Względna pierwszość dowolnych dwóch
czynników i lemat Euklidesa pokazują
indukcyjnie, że \(f_1\cdots f_s\mid g\).
Stąd \([g]=0\); odwzorowanie jest różnowartościowe.

Dla surjektywności połóżmy
\(h_i=\prod_{j\ne i}f_j\).
Ponieważ \(\gcd(h_i,f_i)=1\),
istnieją \(u_i,v_i\in K[t]\) takie, że
\(u_ih_i+v_if_i=1\).
Niech \(e_i=u_ih_i\).
Wtedy \(e_i\equiv1\pmod{f_i}\)
oraz \(e_i\equiv0\pmod{f_j}\)
dla \(j\ne i\).
Dla dowolnych klas \([g_i]_{f_i}\)
wystarczy obrać
\(g=\sum_i g_i e_i\):
jego klasa modulo \(f_j\) jest
dokładnie \([g_j]_{f_j}\).
Zatem \(\Phi\) jest surjektywne.
\end{proof}

\begin{przyklad}[Moduł cykliczny z dwoma czynnikami]
\label{prz:lin20-crt}
Nad \(\mathbb Q\) wielomiany
\(p=t-1\) i \(q=t^2+1\)
są względnie pierwsze.
Zastosowanie Twierdzenia~\ref{thm:lin20-crt}
do \(f=pq\) daje
\[
 \mathbb Q[t]/\bigl((t-1)(t^2+1)\bigr)
 \cong \mathbb Q[t]/(t-1)\oplus
          \mathbb Q[t]/(t^2+1).
\]
Zgodnie ze Stwierdzeniem~\ref{prop:lin19-companion-module}
operator mnożenia przez \(t\)
ma zatem dwie podobne postacie macierzowe:
jeden blok \(C((t-1)(t^2+1))\)
albo dwa bloki \(C(t-1)\oplus C(t^2+1)\).
Nie oznacza to, że ma on dwa
czynniki niezmiennicze:
postać Frobeniusa z
Twierdzenia~\ref{thm:lin19-frobenius}
ma w tym przykładzie \emph{jeden}
czynnik niezmienniczy \(f\).
\end{przyklad}

\section{Istnienie i jednoznaczność bloków pierwotnych}
\label{sec:lin20-primary-blocks}

\begin{twierdzenie}[Postać blokowa przez dzielniki elementarne]
\label{thm:lin20-elementary}
Niech \(T\in\operatorname{End}_K(V)\),
gdzie \(V\) jest skończenie wymiarową
przestrzenią nad ciałem \(K\).
Istnieje baza \(V\), w której macierz
\(T\) jest sumą blokową
\[
 \bigoplus_{p}\ \bigoplus_{j=1}^{r_p}
 C\bigl(p^{a_{p,j}}\bigr).
\]
W tym zapisie \(p\) przebiega skończony
zbiór różnych monicznych wielomianów
nierozkładalnych z \(K[t]\),
\(r_p\ge1\), a
\(1\le a_{p,1}\le\cdots\le a_{p,r_p}\).
Cała rodzina potęg \(p^{a_{p,j}}\),
z uwzględnieniem ich powtórzeń,
jest jednoznacznie określona przez \(T\).
Dla \(V=\{0\}\) zbiór wielomianów \(p\)
i suma blokowa są puste.

Ponadto, jeśli \(n=\dim_K V\),
to
\[
 n=\sum_p\deg(p)\sum_{j=1}^{r_p}a_{p,j},
 \qquad
 \chi_T=\prod_p p^{\,\sum_j a_{p,j}},
 \qquad
 m_T=\prod_p p^{\,a_{p,r_p}}.
\]
Dla pustej rodziny przyjmujemy
pustą sumę równą \(0\), a pusty
iloczyn równy \(1\).
\end{twierdzenie}
\begin{proof}
Niech \(n\ge1\).
Z Twierdzenia~\ref{thm:lin19-frobenius}
operator ma jednoznaczne moniczne
czynniki niezmiennicze
\(f_1\mid\cdots\mid f_r\)
o dodatnich stopniach oraz
\[
 V_T\cong\bigoplus_{i=1}^{r}K[t]/(f_i).
\]
Wybierzmy zbiór \(\mathcal P\)
wszystkich monicznych wielomianów
nierozkładalnych dzielących choć jeden \(f_i\).
Z jednoznaczności rozkładu wielomianów
dla każdego \(p\in\mathcal P\)
istnieją jednoznaczne wykładniki
\(a_{p,i}\ge0\), dla których
\[
 f_i=\prod_{p\in\mathcal P}p^{a_{p,i}}.
\]
Podzielność \(f_i\mid f_{i+1}\)
implikuje
\(a_{p,i}\le a_{p,i+1}\) dla każdego \(p\).
Wielomiany \(p^{a_{p,i}}\)
o dodatnich wykładnikach dla ustalonego \(i\)
są parami względnie pierwsze.
Z Twierdzenia~\ref{thm:lin20-crt}
wynika zatem
\[
 K[t]/(f_i)\cong
 \bigoplus_{\substack{p\in\mathcal P\\a_{p,i}>0}}
 K[t]/\bigl(p^{a_{p,i}}\bigr).
\]
Po podstawieniu tych izomorfizmów
i zebraniu wszystkich składników
dla ustalonego \(p\)
otrzymujemy rozkład modułu \(V_T\)
na sumę modułów cyklicznych
\[
 V_T\cong\bigoplus_{p\in\mathcal P}
        \bigoplus_{j=1}^{r_p}
        K[t]/\bigl(p^{a_{p,j}}\bigr),
\]
gdzie pominięto zerowe wykładniki
i zachowano rosnącą kolejność pozostałych.
Na każdym składniku działanie \(t\)
ma macierz \(C(p^{a_{p,j}})\)
na mocy Stwierdzenia~\ref{prop:lin19-companion-module}.
Izomorfizm modułów przeplata operatory,
więc po połączeniu baz dostajemy
żądaną postać blokową.

Uzasadnijmy jednoznaczność także wtedy,
gdy dane jest \emph{inne} przedstawienie
operatora jako sumy takich bloków.
Dla każdego występującego \(p\)
uporządkujmy jego dodatnie wykładniki
niemalejąco i oznaczmy ich liczbę przez \(r_p\).
Niech \(r=\max_p r_p\).
Uzupełnijmy każdy ciąg
\((a_{p,1},\ldots,a_{p,r_p})\)
zerami \emph{z lewej strony}
do wspólnej długości \(r\),
uzyskując wykładniki \(b_{p,1},\ldots,b_{p,r}\).
Zdefiniujmy wielomiany
\[
 g_i=\prod_p p^{b_{p,i}},
 \qquad i=1,\ldots,r.
\]
Z niemalejącego uporządkowania wykładników
wynika \(g_i\mid g_{i+1}\).
Co najmniej jeden wielomian \(p\)
ma \(r_p=r\), więc wszystkie \(g_i\)
mają dodatnie stopnie.
Ponowne zastosowanie chińskiego
twierdzenia o resztach, tym razem
dla każdego ustalonego \(i\), pokazuje
\[
 \bigoplus_p\bigoplus_{j=1}^{r_p}
 K[t]/\bigl(p^{a_{p,j}}\bigr)
 \cong\bigoplus_{i=1}^{r}K[t]/(g_i).
\]
Prawa strona jest postacią Frobeniusa
o łańcuchu podzielności \(g_1\mid\cdots\mid g_r\).
Z jednoznaczności w
Twierdzeniu~\ref{thm:lin19-frobenius}
wynika, że \(g_i=f_i\) dla wszystkich \(i\).
Jednoznaczny rozkład każdego \(g_i\)
na potęgi nierozkładalne wyznacza
zatem także wyjściową rodzinę
dzielników elementarnych, wraz z krotnościami.

Wymiar modułu \(K[t]/(p^a)\)
jest równy \(a\deg p\),
gdyż klasy potęg \(1,t,\ldots,t^{a\deg p-1}\)
tworzą bazę.
Dodając wymiary składników,
dostajemy wzór na \(n\).
Wielomian charakterystyczny
bloku \(C(p^a)\) jest równy \(p^a\),
więc dla sumy blokowej
\(\chi_T=\prod_p p^{\sum_j a_{p,j}}\).
Wielomian zeruje całą sumę
dokładnie wtedy, gdy zeruje wszystkie bloki;
musi więc być podzielny przez
każdą potęgę \(p^{a_{p,j}}\).
Najmniejsza wspólna wielokrotność
tych wielomianów to
\(\prod_p p^{a_{p,r_p}}\),
czyli \(m_T\).

Jeżeli \(n=0\), wszystkie
wnioski wynikają bezpośrednio
z konwencji o pustych sumach i iloczynach.
\end{proof}

\begin{wniosek}[Czynniki niezmiennicze a dzielniki elementarne]
\label{cor:lin20-reconstruct}
Niech \(T\) działa na niezerowej
skończenie wymiarowej przestrzeni nad \(K\).
Dla każdego monicznego
nierozkładalnego \(p\) występującego
w dzielnikach elementarnych \(T\)
zapiszmy rosnący ciąg dodatnich
wykładników \(a_{p,1},\ldots,a_{p,r_p}\).
Jeżeli \(r=\max_p r_p\),
to po uzupełnieniu tych ciągów
zerami z lewej strony do długości \(r\)
otrzymamy liczby \(b_{p,i}\), a
\[
 f_i=\prod_p p^{b_{p,i}},
 \qquad i=1,\ldots,r,
\]
są dokładnie czynnikami niezmienniczymi \(T\).
\end{wniosek}
\begin{proof}
W drugiej części dowodu
Twierdzenia~\ref{thm:lin20-elementary}
wykazaliśmy, że iloczyny te
mają dodatnie stopnie,
tworzą łańcuch podzielności
i opisują moduł \(V_T\) jako
sumę \(\bigoplus_i K[t]/(f_i)\).
Z jednoznaczności postaci Frobeniusa
wynika teza.
\end{proof}

\begin{wniosek}[Kryterium podobieństwa przez dzielniki elementarne]
\label{cor:lin20-similarity}
Niech \(A,B\in\operatorname{Mat}_{n\times n}(K)\).
Macierze \(A\) i \(B\) są podobne nad \(K\)
wtedy i tylko wtedy, gdy mają
te same dzielniki elementarne
z tymi samymi krotnościami.
\end{wniosek}
\begin{proof}
Podobieństwo zachowuje moduł operatora,
zatem czynniki niezmiennicze
na mocy Wniosku~\ref{cor:lin19-similarity},
a w konsekwencji ich rozkłady
na potęgi wielomianów nierozkładalnych.
Odwrotnie, wspólna rodzina dzielników
elementarnych daje przez
Wniosek~\ref{cor:lin20-reconstruct}
te same czynniki niezmiennicze,
więc macierze są podobne na mocy
Wniosku~\ref{cor:lin19-similarity}.
Dla \(n=0\) obie macierze są puste,
co również spełnia tezę.
\end{proof}

\section{Odczytywanie bloków z wymiarów jąder}
\label{sec:lin20-kernels}

\begin{twierdzenie}[Wymiary jąder potęg wielomianu nierozkładalnego]
\label{thm:lin20-kernel-count}
Niech \(p\in K[t]\) będzie monicznym
wielomianem nierozkładalnym stopnia
\(d\ge1\), a \(T\in\operatorname{End}_K(V)\),
gdzie \(V\) jest skończenie wymiarowe.
Dla \(k\ge0\) połóżmy
\[
 D_{p,k}=\dim_K\ker\bigl(p(T)^k\bigr).
\]
Niech \(a_{p,1},\ldots,a_{p,r_p}\ge1\)
oznaczają wykładniki wszystkich
dzielników elementarnych \(T\)
będących potęgami \(p\), z powtórzeniami.
Jeżeli takich dzielników nie ma,
przyjmujemy \(r_p=0\).
Wówczas dla każdego \(k\ge0\)
\[
 D_{p,k}
 =d\sum_{j=1}^{r_p}\min(k,a_{p,j}).
\]
W szczególności dla \(k\ge1\)
liczba bloków \(C(p^a)\) o \(a\ge k\)
jest równa
\[
 B_{p,k}=\frac{D_{p,k}-D_{p,k-1}}{d},
\]
zaś liczba bloków o \(a=k\) wynosi
\[
 B_{p,k}-B_{p,k+1}
 =\frac{2D_{p,k}-D_{p,k-1}-D_{p,k+1}}{d}.
\]
\end{twierdzenie}
\begin{proof}
Wystarczy obliczyć jądro na każdym
bloku pierwotnym z
Twierdzenia~\ref{thm:lin20-elementary}.
Na module \(R=K[t]/(p^a)\)
operator \(p(T)^k\)
jest mnożeniem przez \(p^k\).

Jeżeli \(0\le k<a\), to
\([p^kg]_{p^a}=0\) dokładnie wtedy,
gdy \(p^a\mid p^kg\), czyli
\(p^{a-k}\mid g\).
Stąd
\[
 \ker p(T)^k
 = (p^{a-k})/(p^a)
 \cong K[t]/(p^k)
\]
jako przestrzenie nad \(K\).
W ostatnim izomorfizmie
klasie \([h]_{p^k}\) przyporządkowuje się
klasę \([p^{a-k}h]_{p^a}\).
Jej wymiar jest równy \(kd\).
Jeśli \(k\ge a\), mnożenie przez \(p^k\)
jest zerowe na całym \(R\),
zatem wymiar jądra wynosi \(ad\).
Oba przypadki łączy formuła
\(d\min(k,a)\); dla \(k=0\)
jądro operatora identycznościowego
jest zerowe.

Na innym bloku
\(K[t]/(q^b)\), gdzie \(q\ne p\)
jest monicznym wielomianem
nierozkładalnym, mamy
\(\gcd(p^k,q^b)=1\).
Tożsamość Bézouta
\(up^k+vq^b=1\)
pokazuje, że mnożenie przez \(p^k\)
ma odwrotność daną przez mnożenie
przez \(u\) modulo \(q^b\).
Zatem jego jądro na tym bloku
jest zerowe.
Ponieważ jądro sumy prostej operatorów
jest sumą prostą ich jąder,
dodanie wkładów wszystkich bloków
daje pierwszy wzór.

Dla ustalonego \(a\ge1\)
różnica \(\min(k,a)-\min(k-1,a)\)
jest równa \(1\), jeśli \(a\ge k\),
i \(0\) w przeciwnym przypadku.
Po odjęciu wzorów dla \(D_{p,k}\)
i \(D_{p,k-1}\) dostajemy
\(B_{p,k}\), liczbę bloków długości
co najmniej \(k\).
Różnica \(B_{p,k}-B_{p,k+1}\)
zlicza dokładnie bloki długości \(k\);
po podstawieniu otrzymujemy
drugi podany wzór.
\end{proof}

\begin{uwaga}[Zakończenie filtracji]
\label{rem:lin20-kernel-stabilization}
Jeśli największy wykładnik
w dzielnikach elementarnych
związanych z \(p\) jest równy \(a\),
to \(D_{p,k}=D_{p,a}\) dla wszystkich
\(k\ge a\).
Wartość \(D_{p,a}\) jest wymiarem
\(p\)-pierwotnego składnika z rozkładu
z Rozdziału~\ref{ch:lin16-niezmiennicze}.
Gdy \(p\) nie występuje
w żadnym dzielniku elementarnym,
wszystkie \(D_{p,k}\) są zerowe.
\end{uwaga}

\section{Związek z postacią Jordana}
\label{sec:lin20-jordan}

\begin{stwierdzenie}[Blok dla wielomianu liniowego]
\label{prop:lin20-linear-primary}
Niech \(\lambda\in K\) i \(a\ge1\).
Blok pierwotny \(C((t-\lambda)^a)\)
jest podobny nad \(K\)
do bloku Jordana \(J_a(\lambda)\)
z Rozdziału~\ref{ch:lin17-jordan}.
\end{stwierdzenie}
\begin{proof}
Na module \(R=K[t]/((t-\lambda)^a)\)
oznaczmy \(u=t-\lambda\).
Klasy
\[
 u^{a-1},u^{a-2},\ldots,u,1
\]
tworzą bazę nad \(K\):
wielomiany \(1,u,\ldots,u^{a-1}\)
rozpinają wszystkie klasy przez dzielenie
z resztą, a niezależność wynika z faktu,
że niezerowy wielomian w \(u\)
stopnia mniejszego niż \(a\)
nie jest podzielny przez \(u^a\).
W tej bazie operator mnożenia przez
\(t=\lambda+u\) spełnia
\[
 t\cdot u^{a-1}=\lambda u^{a-1},
 \qquad
 t\cdot u^{a-j}
   =\lambda u^{a-j}+u^{a-j+1}
       \quad(2\le j\le a).
\]
Jego macierz ma więc na przekątnej
\(\lambda\), jedynki bezpośrednio
nad przekątną i zera poza tym,
czyli jest równa \(J_a(\lambda)\).
Jednocześnie w bazie
\(1,t,\ldots,t^{a-1}\) jest to
\(C((t-\lambda)^a)\)
na mocy Stwierdzenia~\ref{prop:lin19-companion-module}.
Są to zatem macierze tego samego
operatora w dwóch bazach.
\end{proof}

\begin{wniosek}[Dzielniki elementarne i bloki Jordana]
\label{cor:lin20-jordan}
Niech \(T\in\operatorname{End}_K(V)\),
gdzie \(V\) jest skończenie wymiarowe,
a jego wielomian charakterystyczny
rozkłada się nad \(K\) na czynniki liniowe.
Wtedy dzielniki elementarne \(T\)
są dokładnie wielomianami
\((t-\lambda)^a\) odpowiadającymi
blokom Jordana \(J_a(\lambda)\),
z zachowaniem krotności.
\end{wniosek}
\begin{proof}
Każdy nierozkładalny czynnik
wielomianu charakterystycznego
jest w tym przypadku liniowy,
a każdy dzielnik elementarny
jest potęgą takiego czynnika.
Twierdzenie~\ref{thm:lin20-elementary}
rozkłada operator na bloki pierwotne,
które na mocy
Stwierdzenia~\ref{prop:lin20-linear-primary}
są podobne do odpowiednich
bloków Jordana.
Jednoznaczność obu klasyfikacji
identyfikuje także ich krotności.
\end{proof}

\section{Przykłady i obliczenia}
\label{sec:lin20-examples}

\begin{przyklad}[Dwa bloki nilpotentne]
\label{prz:lin20-nilpotent}
Niech \(A=J_3(0)\oplus J_1(0)\)
nad dowolnym ciałem \(K\).
Dzielniki elementarne to \(t^3\) oraz \(t\).
Dla \(p=t\) otrzymujemy kolejno
\[
 D_{t,0}=0,\quad D_{t,1}=2,\quad
 D_{t,2}=3,\quad D_{t,3}=4,\quad
 D_{t,4}=4.
\]
Przyrosty \(2,1,1,0\)
oznaczają: dwa bloki długości
co najmniej \(1\), jeden długości
co najmniej \(2\), jeden długości
co najmniej \(3\), żaden długości
co najmniej \(4\).
Zatem dokładnie jeden blok
ma długość \(1\), a jeden długość \(3\).
Reguła z Wniosku~\ref{cor:lin20-reconstruct}
daje czynniki niezmiennicze
\(f_1=t\) i \(f_2=t^3\);
stąd \(\chi_A=t^4\) oraz \(m_A=t^3\).
\end{przyklad}

\begin{przyklad}[Dwie wartości własne i wyrównanie wykładników]
\label{prz:lin20-diagonal}
Nad \(\mathbb Q\) rozważmy
\(A=\operatorname{diag}(1,1,2)\).
Wszystkie bloki są jednowymiarowe,
a dzielniki elementarne to
\[
 p=t-1,\quad p=t-1,\quad q=t-2.
\]
Wykładniki dla \(p\) wynoszą \((1,1)\),
natomiast dla \(q\) jest jeden
wykładnik \((1)\).
Uzupełniając drugi ciąg z lewej
strony do długości dwóch,
otrzymujemy odpowiednio
\((1,1)\) i \((0,1)\).
Iloczyny na odpowiadających
pozycjach dają
\[
 f_1=p,\qquad f_2=pq.
\]
Są to dokładnie czynniki niezmiennicze
obliczone w
Przykładzie~\ref{prz:lin19-diagonal}.
Zauważmy, że \emph{nie wolno}
przyjąć \(p^2\) i \(q\) za dwa
czynniki niezmiennicze: choć ich
iloczyn jest równy \(\chi_A=p^2q\),
wielomiany nie tworzą łańcucha
podzielności. Ponadto ich najmniejsza
wspólna wielokrotność wynosiłaby \(p^2q\),
podczas gdy \(m_A=pq\).
\end{przyklad}

\begin{przyklad}[Wielomian nierozkładalny stopnia dwa]
\label{prz:lin20-real}
Niech \(K=\mathbb R\),
\(q=t^2+1\) oraz
\[
 A=C(q)\oplus C(q^2),
 \qquad
 C(q^2)=
 \begin{pmatrix}
 0&0&0&-1\\
 1&0&0&0\\
 0&1&0&-2\\
 0&0&1&0
 \end{pmatrix}.
\]
Ponieważ \(q^2=t^4+2t^2+1\),
ostatnia kolumna tej macierzy
jest zgodna z definicją
\(\,C(q^2)\) z
Rozdziału~\ref{ch:lin19-frobenius}.
Dzielniki elementarne \(A\)
to \(q,q^2\); czynniki niezmiennicze
są równe \(f_1=q,f_2=q^2\),
zaś \(\chi_A=q^3\), \(m_A=q^2\).

Ponieważ \(\deg q=2\),
Twierdzenie~\ref{thm:lin20-kernel-count}
daje
\[
 D_{q,0}=0,\quad D_{q,1}=2(1+1)=4,
 \quad D_{q,2}=2(1+2)=6.
\]
Zatem \(B_{q,1}=2\), \(B_{q,2}=1\),
\(B_{q,3}=0\): są dokładnie dwa
bloki \(q\)-pierwotne, z których jeden
ma wykładnik \(1\), a drugi \(2\).
Wielomian \(q\) nie ma rzeczywistych
pierwiastków, więc macierz \(A\)
nie ma rzeczywistych wartości własnych.
Nie przeszkadza to w opisanej klasyfikacji.
\end{przyklad}

\begin{przyklad}[Blok pierwotny nad ciałem dwuelementowym]
\label{prz:lin20-finite}
Niech \(K=\mathbb F_2\) oraz
\(p=t^2+t+1\).
Wielomian \(p\) jest nierozkładalny:
ma stopień \(2\) i nie zeruje się
ani dla \(t=0\), ani dla \(t=1\).
Ponieważ w charakterystyce \(2\)
zachodzi \(p^2=t^4+t^2+1\),
macierz
\[
 A=C(p^2)=
 \begin{pmatrix}
 0&0&0&1\\
 1&0&0&0\\
 0&1&0&1\\
 0&0&1&0
 \end{pmatrix}
\]
ma cztery wiersze i kolumny.
Jej jedynym dzielnikiem elementarnym
i jedynym czynnikiem niezmienniczym
jest \(p^2\), więc
\(\chi_A=m_A=p^2\).
Filtracja jąder daje
\(D_{p,0}=0\), \(D_{p,1}=2\),
\(D_{p,2}=4\) oraz
\(D_{p,k}=4\) dla \(k\ge2\).
Odczytujemy z niej jeden blok
o wykładniku \(2\).
Macierz nie ma wartości własnych
w \(\mathbb F_2\), mimo że
ma kompletną postać pierwotną.
\end{przyklad}

\begin{przyklad}[Wielomiany charakterystyczny i minimalny nie wystarczają]
\label{prz:lin20-not-enough}
Niech nad dowolnym ciałem \(K\)
\[
 A=J_3(0)\oplus J_3(0),\qquad
 B=J_3(0)\oplus J_2(0)\oplus J_1(0).
\]
Obie macierze mają rozmiar \(6\)
oraz \(\chi_A=\chi_B=t^6\) i
\(m_A=m_B=t^3\).
Nie są jednak podobne:
\(\dim\ker A=2\), natomiast
\(\dim\ker B=3\), ponieważ
każdy blok Jordana z wartością
własną \(0\) wnosi jeden wymiar
do jądra. Ich dzielniki elementarne
to odpowiednio \((t^3,t^3)\)
oraz \((t,t^2,t^3)\).
Przykład ilustruje, dlaczego
sama znajomość \(\chi_T\) i \(m_T\)
nie wystarcza do klasyfikacji.
\end{przyklad}

\section{Podsumowanie}
\label{sec:lin20-summary}

Dzielniki elementarne są potęgami monicznych
wielomianów nierozkładalnych
występujących w czynnikach niezmienniczych.
Chińskie twierdzenie o resztach
pozwala zamienić każdy moduł cykliczny
\(K[t]/(f)\) na sumę modułów
\(K[t]/(p^a)\).
Otrzymana suma bloków pierwotnych
jest jednoznaczna z dokładnością
do ich kolejności.

Obie klasyfikacje są równoważne:
z dzielników elementarnych odzyskujemy
czynniki niezmiennicze przez uporządkowanie
wykładników i dopełnianie zerami z lewej.
Różnice wymiarów jąder \(p(T)^k\)
zliczają bloki o kolejnych długościach.
Jeśli wszystkie nierozkładalne czynniki
wielomianu charakterystycznego są liniowe,
bloki pierwotne odpowiadają
dokładnie blokom Jordana.
